\documentclass[10pt,a4paper]{amsart}
\usepackage[margin=19mm]{geometry}
\usepackage{amsmath,amssymb,mathtools}
\usepackage[scr=rsfs,cal=euler]{mathalfa}
\usepackage{xfrac}
\usepackage[unicode,hidelinks,bookmarks=false]{hyperref}
\newcommand{\ie}{i.e.\ }
\newcommand{\cf}{cf.\ }
\newcommand{\resp}{respectively\ }
\newcommand{\Subsection}{\subsection}
\providecommand{\nobreakdash}{\nobreak\hbox{-}}
\setlength{\emergencystretch}{2em}
\multlinegap0pt
\usepackage{amssymb}
\usepackage{mathtools}
\usepackage{xcolor}
\newtheorem{proposition}{Proposition}
\newcommand{\F}{\mathcal F}
\def\R{{\mathbb R}}
\newcommand{\abs}[1]{\left\lvert#1\right\rvert}
\newcommand{\angles}[1]{\left\langle#1\right\rangle}
\newcommand{\norm}[1]{\lVert#1\rVert}
\newcommand{\set}[1]{\left\{#1\right\}}
\newcommand{\sgn}{\operatorname{sgn}}
\title{A Coarse-Lipschitz Embedding of $c_0$\\ into a Separable Dual Banach Space}
\author{Bunyamin Sari}
\date{Source: arXiv:2608.04117v2 (CC BY 4.0)}
\begin{document}
\maketitle

\noindent\emph{Source and licence.} A Coarse-Lipschitz Embedding of $c_0$\\ into a Separable Dual Banach Space. Version arXiv:2608.04117v2 (CC BY 4.0), distributed by arXiv under the Creative Commons Attribution 4.0 International licence, http://creativecommons.org/licenses/by/4.0/, as recorded in the arXiv licence metadata for this exact version and shown on the abstract page https://arxiv.org/abs/2608.04117v2. Full licence notice: \texttt{licenses/arxiv-2608.04117-CC-BY-4.0.txt}.\par\smallskip

\noindent\textit{Editorial scope.} Counted unit: the article's proposition prop:parent-edge-basis (source0.tex lines 323--337). The article's complete original proof of it is delivered with it (source0.tex lines 340--403, one \texttt{proof} environment of 64 lines). No mathematics is written, repaired or re-derived: the statement and the proof are the article's own lines, in the article's own notation, and no retained line is substituted or re-flowed.  Retained uncounted as the source-local context the proof needs: lines 244--246.  Nothing was omitted from any retained span, so the package carries no omission record.  No retained line contains a citation, so the wrapper carries no bibliography and no bibliography entry.  No retained line contains a figure environment, an image-inclusion command or drawing code, so no image is delivered and the package declares no asset dependency.  Editorial formatting only, in addition: an AMS \texttt{amsart} preamble that restates the article's own declarations and macros used by the retained lines, namely the environments \texttt{proposition} on separate counters, together with the standard packages those lines need.  The retained environments print in this document as Proposition 1 in order of appearance, whereas the article prints the same items with the numbering of its own full document.  The complete native source of the article as distributed in the arXiv e-print for this exact version is bundled unchanged as \texttt{sources/206804/source0.tex}.
\begin{equation}\label{eq:enumeration}
       G_R=\set{0,x_1,x_2,\ldots}
\end{equation}

\begin{proposition}\label{prop:parent-edge-basis}
With respect to the enumeration \eqref{eq:enumeration}, the edge vectors (molecules)
$$
b_n=\delta(x_n)-\delta(px_n),
\qquad n\geq1,
$$
form a normalized Schauder basis of $\F(G_R)$ with basis constant at most two.  When $R=1$,
the basis is monotone.  Moreover, for every finite scalar sequence
$(a_j)$, we have
\begin{equation}\label{eq:l1-equivalence}
\frac{1}{2R}\sum_j\abs{a_j}
\leq\Big\|\sum_j a_jb_j\Big\|
\leq\sum_j\abs{a_j}.
\end{equation}
\end{proposition}

\begin{proof}
Since $px_j\in A_{j-1}$, $b_j$ is the molecule associated with the parent edge $(px_j, x_j)$. That is, it is the new edge introduced when $x_j$ is added the to $A_{j-1}$ (this is the reason for the formula $b_j=\delta(x_j)-\delta(px_j)$). We check that 
$$
P_nb_j=
\begin{cases}
b_j,&j\leq n,\\
0,&j>n.
\end{cases}
$$
If $j\leq n$, then $r_n$ fixes both $x_j$ and $px_j$.  Suppose that $j>n$,
and let
$$
k=\min\set{l\geq1:p^lx_j\in A_n}.
$$
Then $r_nx_j=p^kx_j$.  The first point on the parent chain of $px_j$ which
belongs to $A_n$ is the same point, since
$$
p^{k-1}(px_j)=p^kx_j.
$$
Thus $r_n(px_j)=r_nx_j$, and consequently $P_nb_j=0$. Thus $P_n$'s are $n$th partial sum operators, and recall that $\norm{P_n}\le 2$.

Note that
$$
\delta(x_j)=\sum_{l=0}^{\norm{x_j}_\infty-1}
\bigl(\delta(p^lx_j)-\delta(p^{l+1}x_j)\bigr),
$$
and every summand on the right is one of the vectors $b_i$ with $i\leq j$.
Therefore
$$
\operatorname{span}\set{b_1,\ldots,b_n}
=\operatorname{span}\set{\delta(x_1),\ldots,\delta(x_n)}
$$
for every $n$.  In particular, $(b_j)$ is linearly independent and has dense
linear span in $\F(G_R)$.

It remains to prove \eqref{eq:l1-equivalence}.  Let $(a_j)$ be finitely
supported and write
$$
\mu=\sum_j a_jb_j=\sum_k c_k\delta(x_k).
$$
Since $G_R$ is $1$-separated, the function $f:G_R\to\R$ defined by
$$
f(0)=0,
\qquad
f(x_k)=\frac12\sgn(c_k)
$$
is $1$-Lipschitz.  Hence
$$
\norm{\mu}\geq\angles{\mu,f}=\frac12\sum_k\abs{c_k}.
$$
On the other hand, each Dirac vector $\delta(x_k)$ is the sum of the
parent-edge molecules along its path from $0$ to $x_k$, and this path has
length at most $R$.  Expanding each $\delta(x_k)$ in this way and comparing
the unique $(b_j)$-expansions of $\mu$ gives
$$
\sum_j\abs{a_j}\leq R\sum_k\abs{c_k}\leq2R\norm{\mu}.
$$
Since $\norm{b_j}=d(x_j,px_j)=1$, we also have
$$
\norm{\mu}\leq\sum_j\abs{a_j}.
$$
This proves \eqref{eq:l1-equivalence}.

\end{proof}

\end{document}
